\section{Heliocentrism Redux}

The point here is not to rekindle the dreary ``debate'' between heliocentrism and geocentrism, but rather to use it to elucidate the nature of scientific proof. So this cannot an exhaustive explanation.

First of all, let's clear up a little detail. The theory is not that the earth revolves around the Sun, Rather, the Sun and the Earth both travel around a focus of an ellipse. The focus lies at the barycenter, i.e., the common center of gravity of the earth (actually the barycenter of the entire solar system).

\begin{quotex}
The Sun is not at the geometric center of any planet's orbit, but rather approximately at one focus of the elliptical orbit. Furthermore, to the extent that a planet's mass cannot be neglected in comparison to the Sun's mass, the center of gravity of the Solar System is displaced slightly away from the center of the Sun. 
\flright{\textit{Wikipedia}}
\end{quotex}


\paragraph{Abductive Logic}


\begin{quotex}
Abductive reasoning is a form of logical inference. It starts with a set of observations and then seeks the simplest and most likely conclusion from the observations. This process, unlike deductive reasoning, yields a plausible conclusion but does not positively verify it.
\end{quotex}


Abductive logic is the ``scientific method''. The method begins with a set of observations, also called the ``appearances'' or empirical data. Then a hypothesis is found that would account for, or explain, the appearances. A better hypothesis would be able to predict future appearances. This is the opposite of deductive logic whose conclusions necessarily follow from the premises. In abductive logic that is not the case because the hypothesis may be replaced by a different, or perhaps a better, one.

Abductive logic was more traditionally called ``\textbf{Saving the Appearances}''.

\paragraph{The Religious View}


Geocentrism is often accepted purportedly on religious grounds. It seems to be compatible with Scriptures and seems to give the human race a special place in the universe. However, is it not a dogma of the Catholic church (I don't know about other denominations).

Nevertheless, it is spiritually true (leaving aside the physical question). The Traditional view is not that the Earth is special, but rather it is a remote outpost furthest from Heaven. It is sublunar and closest to hell.

Dante illustrates the metaphysical understanding. The geocentric model represents a spiritual journey that starts with the descent into hell. Then follows a period of purification, then an ascent through the planetary spheres, the stars, and eventually Heaven.

\paragraph{Astronomical Theories}





Careful observations led to three different hypotheses a few hundred years ago, to wit:

\begin{enumerate}


\item \textbf{Ptolemaic}: Geocentric, the earth is the center of the solar system all planets plus the sun revolve around the earth although in complex patterns called epicycles

\item \textbf{Copernican}: Heliocentric, all planets including the Earth revolve around the sun

\item \textbf{Tychonic}: Geocentric, Earth is still, the moon and sun revolve around the Earth, the other planets revolve around the sun
\end{enumerate}


At the time, all three options saved the appearances, that is, they were all able to predict the future positions of the planets. Moreover, the Copernican and Tychonic models are mathematically equivalent:

\begin{quotex}
At the same time, the motions of the planets are mathematically equivalent to the motions in Copernicus' heliocentric system under a simple coordinate transformation, so that, as long as no force law is postulated to explain why the planets move as described, there is no mathematical reason to prefer either the Tychonic or the Copernican system.
\end{quotex}


This raises two interesting questions. How can an ``incorrect'' hypothesis still save the appearances? You might think that bad science leads to bad results, but not necessarily.

The second question is how can the same mathematical model be interpreted in two different ways? After all, there is no standpoint for the astronomer to objectively view the Solar System to see what is ``really going on''.

Obviously, the astronomer will prefer the simpler mathematical model but that does not translate to some imaginary objective viewpoint.

\paragraph{General Relativity}


General relativity supersedes all the above options. In this theory, there is no ``force of gravity''; rather the Earth moves in straight line. However, space is distorted by the mass of the Sun so it may not be a straight line in Euclidean space. The question of whether the Sun moves or the Earth moves is moot. Einstein explains:

\begin{quotex}
Can we formulate physical laws so that they are valid for all CS (= coordinate systems)? If this can be done, our difficulties will be over. We shall then be able to apply the laws of nature to any CS. The struggle, so violent in the early days of science, between the views of Ptolemy and Copernicus would then be quite meaningless. Either CS could be used with equal justification. The two sentences, `the sun is at rest and the Earth moves', or `the sun moves and the Earth is at rest', would simply mean two different conventions concerning two different CS.
\flright{\textit{Evolution of Physics}\footnote{\url{https://www.gornahoor.net/library/EvolutionOfPhysics.pdf}}}
\end{quotex}


In other words, physics by itself cannot determine if the earth moves or the sun moves. That can only be determined in some other, unspecified, way.

We conclude with two principles, incompatible with each other, to put things into perspective.

\paragraph{Copernican Principle}


This principle states that an observer on the Earth does not have a privileged position in the universe because the universe is homogeneous everywhere. The corollary is that the Earth is insignificant in the universe, just one planet among many, given the vast size and age of the universe. But ``vast'' is a relative term; maybe vast in terms of a human life, but there is no other standard to use.

Unfortunately, there is no way to prove the Copernican Principle. Moreover, things may not be so simple. Planets like the earth may be rare. Life may be rarer. Consciousness even rarer, and thought rarer still. The astronomer cannot tell us how beings arose who can create theories about the universe, as if it were self-reflective.

\paragraph{Anthropic Principle}


There are two versions of the anthropic principle.

The weak anthropic principle states that the size and age of the universe are necessary for a planet with life, consciousness, and thought to arise. Apparently, it takes quite a while to form a planet with the correct chemical conditions to support life. Were the size of the universe too small, then it would have collapsed upon itself before such a planet to arise. The point is that the size, age, and physical constants are such as they other for life to arise. This principle does not require a designer.

The strong anthropic principle states that the universe is such as it is precisely so the conscious and sapient life could emerge. Since we know that conscious, sapiential life has emerged, then there must be a physical theory to explain how. But there is no such theory, nor is one forthcoming.

Some things have no physical explanations. As Einstein pointed out, there is no way to determine what is actually in motion. Moreover, in all the equations of physics, time itself has no direction; the equations hold whether or not time is moving forward or back. But somehow we can tell.


\flrightit{Posted on 2021-11-07 by Cologero}

\begin{center}* * *\end{center}

\begin{footnotesize}
\begin{sffamily}
\texttt{veritycacciatrice on 2021-11-08 at 05:47 said:}

The Hermetic metaphor appears everywhere in everything, and I've often wondered when and where this fine art originated.

It is an ongoing conversation/running joke in our house a bit like the identity of the Sea Peoples (currently, they were state sponsored mercenaries).

Egypt seems clear enough as the birth place of the Hermetica, but then one wonders how/why it came about in Egypt.

The Electric Universe Theory, spearheaded by Wal Thornhill and David Talbot have, over the last 40 or 50 years, tied all the ratty ends of the historic helio+geocentric gravity-based debates in to an elegant answer.

Their ideas use Einstein's relativity, Velikovsky's myth analysis and plasma physics/cosmology to explain not only the origins of the myths and religions, but to my mind the possibility of the origins of the Hermetica itself.

Gravity plays so little part as to be almost irrelevant as it is all based on electrically-based plasma physics which can be reproduced in a lab setting. Both theory and practice of this neat answer to our solar system was inspired by cosmic myths of the world using Velikovsky's book as starting point.

To my mind, the reality of it it wipes all other forms of cosmology off the blackboard.

Not everyone is/was impressed by Velikovsky. He is still considered a crank for various reasons, by various parties. But generally it is because they haven't actually sat down and read his work, which very possibly means having to change their entire mental paradigm.

There is absolutely room for God within the shifted paradigm. Spirit is an integral part, and curiously they developed a lab experiment to prove it.

But the most interesting thing to me is it *does* nod to the origins of the Hermetica, the where, how and why, and even perhaps why the Sea Peoples made a quick trip to take what they could just prior to the chaos of the Bronze Age collapse. Good old astrologers and their wacky disciplines.

``The sensational claims of Dr. Immanuel Velikovsky fail to interest me as much as they should, .. because his conclusions were pretty obviously based on incompetent data. .. If in historical times there have been these changes in the structure of the solar system, .. then the laws of Newton are false. .. if Dr. Velikovsky is right, the rest of us are crazy.''

-Harlow Shapley, an astronomer and critic of Velikovsky (who had no jaw to speak of).

\end{sffamily}
\end{footnotesize}

